\documentclass[11pt]{article}
\usepackage[T1]{fontenc}
\usepackage[utf8]{inputenc}
\usepackage{amsmath,amssymb}
\usepackage[margin=1in]{geometry}
\usepackage{hyperref}
\hypersetup{colorlinks=true,urlcolor=blue,linkcolor=blue}
\setlength{\parindent}{0pt}
\setlength{\parskip}{6pt}
\setlength{\emergencystretch}{3em}
\title{P16.1: A Proved Curve Sector and the Remaining Singular-Variety Problem}
\author{Unverified research draft}
\date{4 October 2026}
\begin{document}
\maketitle
This is an unverified mathematical draft. It proves the curve theorem stated below from exact published background inputs and the descent arguments given here. It does not solve the full original atom P16.1. The general unibranch envelope replacement in higher dimension, and the singular finite-energy equation under the discretely valued hypotheses, remain gaps. Hashes, executions and TeX compilation are reproducibility evidence, not mathematical verification.

\section*{1. Exact result and normalization}
Theorem (partial result). Let $k$ be an algebraically closed field of characteristic zero with the trivial absolute value, let $X$ be an integral projective unibranch $k$-curve, and let $L$ be an ample line bundle on $X$. Write $V=\deg_X L>0$, $\theta=c_1(L)$, and $\nu:Y\longrightarrow X$ for the normalization. No smoothness of $X$, normality of $X$, reducedness of the conductor subscheme, or restriction to a skeleton is assumed. Integrality includes reducedness and connectedness. The field is complete because its valuation is trivial. The polarization remains the actual line bundle on $X$.

Use the canonical trivial metric on $L$ as the reference metric. Define $\operatorname{PSH}(X,L)$ using decreasing limits of semipositive model metrics, equivalently the ample-class Fubini--Study construction of Boucksom--Jonsson [BJ]. Define the normalized intersection energy $E_X$ with $E_X(0)=0$ and $E_X(\varphi+c)=E_X(\varphi)+c$, and put $\mathcal E^1(X,L)=\{\varphi\in\operatorname{PSH}(X,L):E_X(\varphi)>-\infty\}$.

There is a unique decreasing-continuous extension of the genuine model Monge--Ampere measure

\[
 \mathrm{MA}^{\rm raw}_{X,L}:\mathcal E^1(X,L)\longrightarrow
 \{\hbox{positive Radon measures on }X^{\rm an}\hbox{ of mass }V\}.
\]

For a positive Radon measure $\mu$, the equation $\mathrm{MA}^{\rm raw}_{X,L}(\varphi)=\mu$ has a solution in $\mathcal E^1(X,L)$ if and only if

\[
 \mu(X^{\rm an})=V,\qquad
 E_X^*(\mu/V):=\sup_{\varphi\in\mathcal E^1(X,L)}
 \left(E_X(\varphi)-V^{-1}\int\varphi\,d\mu\right)<\infty.
\]

The supremum can equivalently be taken over $H_{\rm dom}(X,L)$, the ample model potentials on models dominating the trivial test configuration. The solution is unique modulo one additive real constant. It is uniquely fixed by $\varphi(o)=\sup\varphi=0$, where $o$ is the trivial valuation of $k(X)$. Its metric is unique modulo multiplication by a positive constant. The operator agrees with the intersection-theoretic Monge--Ampere operator on actual $X$, is continuous along decreasing nets in finite energy, and is strongly continuous with the [BJ] strong topologies. Compactness, envelopes, energy differentiation and comparison are justified in Section 7.

This sector changes the field endpoint to the established trivial-valuation endpoint and restricts dimension to one. It is not asserted to extend the original discretely valued theorem in arbitrary dimension. It is a genuine singular-curve theorem with all normalization descents proved, rather than a theorem merely on $Y$.

\section*{2. Exact published inputs and their limits}
[BJ] means Boucksom--Jonsson, *Global pluripotential theory over a trivially valued field*, the supplied arXiv:1801.08229v4. Its standing field is algebraically closed and trivially valued. Theorem A extends the finite-energy operator to possibly singular projective varieties and gives an embedding with dense image; its unconditional surjectivity statement assumes smoothness and characteristic zero or dimension at most two. The body handles reducible varieties, but all positive assertions here are for integral varieties. In this version, Theorem 5.11 equates envelope regularity and weak compactness; Theorem 12.8 equates envelope regularity and finite-energy surjectivity; Conjecture 5.14 proposes envelopes for unibranch varieties. Lemma 5.13 assumes envelope properties on the target. Using that lemma to prove the same envelope property by normalization would be circular.

The background inputs used in the proof are the following exact results. Example 2.5 makes $\max\{\log|\mathfrak a|,-r\}$ a PL function for any ideal on any projective variety and rational $r>0$. Example 2.6 gives the ray generators of PL on a smooth curve. Proposition 3.6(vi) reflects model psh positivity under a surjective projective morphism, for a function already in PL on the target. Proposition 3.10 uniformly approximates a nef PL potential in a fixed ample class by $H_{\rm dom}$ potentials. Theorem 4.15(iii) gives decreasing approximation in that fixed ample class; Theorem 4.7 supplies decreasing closure. Example 4.13 is the smooth-curve convex-ray criterion. Theorem 5.20 gives the envelope property on smooth connected varieties for nef classes when characteristic is zero or dimension is at most two. All its hypotheses hold on $Y$ here.

[BJ] Theorems 7.1 and 7.18 and Proposition 7.19 construct energy and genuine finite-energy MA on singular projective varieties, with normalization identities (7.5) and (7.18). Example 7.20 gives the curve Laplacian. Theorems 7.40 and 7.44 give locality and comparison without smoothness. Theorem 8.5 and Corollary 8.6 differentiate the downward extension of energy on singular varieties; identifying it with energy of a psh envelope requires the envelope property. Proposition 12.7 supplies strong continuity. These results are published background, not new claims proved solely by running a program. Their precise application checks and selected literal proofs are retained in source-extracts/source-provenance.json.

The cached arXiv:1201.0188 is its actual supplied v1. Its PDF watermark says 30 December 2011, while its internal date says 29 October 2018; both are retained literally, without relabelling the cache. Theorem A assumes a complete discretely valued field of residue characteristic zero, a smooth projective variety admitting descent to a one-variable function field whose completion is the ground field, ample $L$, raw mass $\deg_L(X)$, and support on the dual complex of an SNC model. The polarization need not descend to the function field. This is not an arbitrary-support finite-energy theorem. The supplied 1201.0187v2 proves compactness and regularization under smooth projective, discrete-valuation and residue-characteristic-zero hypotheses.

The cached 1805.01016v3 is Boucksom--Eriksson, *Spaces of norms, determinant of cohomology and Fekete points in non-Archimedean geometry*. Theorem A concerns continuous psh metrics on geometrically reduced projective schemes over complete valued fields. Corollary B retains continuity of envelopes as a hypothesis. Proposition 9.14 is a continuous-metric energy identity, not a general finite-energy existence theorem. The synthetic 2307.01697v1, Theorem A, completes the space of finite-energy measures and gives density of test-MA measures; it does not construct every such measure as MA of an actual singular-space finite-energy potential. Pille--Schneider, 2209.04879v2, Remark 4.4 distinguishes the arbitrary-field definition of finite energy from the unknown general extension of its MA operator. The September 2025 Gubler--Rabinoff 2506.13548v2 explicitly says that regularization and defining MA are not the subject of that paper. Its singular-curve subharmonic theory is not silently substituted for the global equation theory. The other three supplied sources are background caches, not indispensable inputs to the new theorem.

\section*{3. Analytic normalization, conductor and PL descent}
Normalization is finite for a variety of finite type over a field. A normal curve over the perfect field $k$ is smooth. Thus $Y$ is a smooth integral projective curve. A unibranch curve has exactly one point $q$ above each closed point $p$ of $X$; since $k$ is algebraically closed, both residue fields are $k$. For this curve argument these are the precise consequences of unibranchness that are needed.

Every point of $Y^{\rm an}$ is either the trivial valuation $o$, a point $t\operatorname{ord}_q$ for $q\in Y(k)$ and $0<t<\infty$, or the endpoint $q_\infty$ with support $q$. Indeed a valuation of $k(Y)$ has a center by properness. If that center is the generic point its valuation is trivial. Otherwise the local ring is a DVR, every unit has value zero, and the valuation is a positive multiple of its order valuation. A semivaluation with nongeneric support is the trivial valuation on the closed point's residue field. The same function field is $k(X)$. A finite birational normalization therefore bijects generic-support semivaluations, and the unique-preimage condition bijects the closed-support semivaluations. Consequently $\nu^{\rm an}$ is bijective. It is continuous, with compact source and Hausdorff target, and is therefore a homeomorphism. This proof does not claim that a normalization of an arbitrary singular variety is an analytic homeomorphism.

For each $p\in X(k)$ let $q$ be its unique preimage. Define the positive integer $e_p$ by

\[
 \mathfrak m_p\mathcal O_{Y,q}=\mathfrak m_q^{e_p}.
\]

The ideal is nonzero and vanishes at $q$, so $1\leq e_p<\infty$. Away from $q$ its extension is the unit ideal. Its canonical ideal norm therefore satisfies

\[
 \nu^*\log|\mathfrak m_p|(t\operatorname{ord}_q)=-e_pt,
 \quad \nu^*\log|\mathfrak m_p|(o)=0,
 \quad \nu^*\log|\mathfrak m_p|(q_\infty)=-\infty,
\]

and is zero on every other ray, including its endpoint. For rational $a>0$ set

\[
 h_{p,a}=e_p^{-1}\max\{\log|\mathfrak m_p|,-e_pa\}\in\operatorname{PL}(X).
\]

Example 2.5 applies to the coherent point ideal on singular $X$. Its pullback is the function $-\min\{t,a\}$ on the $q$-ray and zero on all other rays. These functions, with rational constants, span $\operatorname{PL}(Y)$ by Example 2.6. Pullback always maps flag-ideal functions on $X$ to flag-ideal functions on $Y$. Since $\nu^{\rm an}$ is surjective, pullback is injective. We have proved

\[
 \nu^*:\operatorname{PL}(X)\simeq\operatorname{PL}(Y),
 \qquad \nu^*:\operatorname{PL}(X)_{\mathbb R}\simeq\operatorname{PL}(Y)_{\mathbb R}.
\]

The conductor is $\mathfrak c=\operatorname{Ann}_{\mathcal O_X}(\nu_*\mathcal O_Y/\mathcal O_X)$. It has finite support and may be nonreduced. The proof did not assert that arbitrary sections on $Y$ descend, that powers of such sections descend, or that normalization changes no local algebra. It used actual ideals on $X$. The integer $e_p$ is the point-ideal order, not a conductor exponent. One branch at every conductor point is indispensable for the displayed isolated-ray formula. There are no extra jet conditions on $\mathfrak c$ in this theorem.

\section*{4. Polarization, nef reflection and exact psh descent}
Put $A=\nu^*L$. It is ample: on an integral projective curve ampleness is equivalent to positive degree, and degree is unchanged here. To verify the degree rather than assume it, the finite birational map gives an exact sequence with finite-length quotient $T$,

\[
 0\longrightarrow\mathcal O_X\longrightarrow\nu_*\mathcal O_Y
 \longrightarrow T\longrightarrow0.
\]

Tensor by $L^m$ and use finite pushforward and the projection formula. The difference $\chi(Y,A^m)-\chi(X,L^m)$ is the constant length of $T$. Comparing the linear coefficients of these Hilbert polynomials gives $\deg_Y A=\deg_X L=V$. In particular no degree factor and no normalization of a restriction of $\mu$ occur. Pullback of the canonical metric of $L$ is the canonical metric of $A$, since a local trivialization at the center pulls back to a local trivialization at its inverse image.

Here is the application of nef reflection, with its essential prerequisite explicit. Given $f\in\operatorname{PL}(X)_{\mathbb R}$, choose a test configuration of actual $X$ dominating $X\times\mathbb A^1$ and a vertical divisor determining $f$. Taking the closure of the graph of the generic morphism from $Y\times\mathbb G_m$, and dominating also the trivial configuration of $Y$, gives a projective surjective morphism of test configurations. No smooth resolution or ample pullback along a resolution is needed. Pulling back the representing class preserves nefness. Conversely, if its pullback is nef, each vertical irreducible curve downstairs is dominated by a projective curve upstairs; one obtains such a curve by hyperplane cuts in a component dominating the downstairs curve. The degree of the pulled-back line bundle is the positive degree of that map times the downstairs degree. Thus nonnegative degrees upstairs force nonnegative degrees downstairs. Passing to a common dominating model does not change the model-psh property. This is precisely the argument behind [BJ, Proposition 3.6(vi)]. Therefore

\[
 f\in\operatorname{PL}(X)_{\mathbb R}\text{ is }L\text{-psh}
 \quad\Longleftrightarrow\quad
 \nu^*f\text{ is }A\text{-psh}.
\]

This conclusion cannot be applied to an arbitrary function on $Y$ before its PL descent is proved. Section 3 supplies that missing prerequisite in this curve case.

Now let $\psi\in\operatorname{PSH}(Y,A)$. Approximate it decreasingly by $H_{\rm dom}(Y,A)$, using Theorem 4.15(iii). Descend every approximant through the proved PL isomorphism. Nef reflection makes each descended approximant $L$-psh on actual $X$. The analytic homeomorphism makes them a decreasing family with pointwise limit $\varphi=\psi\circ(\nu^{\rm an})^{-1}$. The limit is usc and generically finite, and decreasing closure makes it $L$-psh. The forward inclusion is ordinary pullback of psh functions, obtained either from Fubini--Study pullback or the same approximation. Hence

\[
 \nu^*:\operatorname{PSH}(X,L)\simeq\operatorname{PSH}(Y,A).
\]

It is a homeomorphism for the weak topology of pointwise convergence on divisorial valuations: those valuations of the common function field are identified. This is a proof of descent before invoking envelopes or compactness, so no circular assumption on singular $X$ has entered. Proposition 3.10 shows further that the descended model approximants can themselves be uniformly approximated in the fixed class by $H_{\rm dom}(X,L)$.

\section*{5. Energy and the genuine operator on actual X}
For $\varphi\in\operatorname{PSH}(X,L)$ write $a=\varphi(o)$ and identify the rays through the proved homeomorphism. Put $b_q(t)=-\frac{d}{dt}\varphi(t\operatorname{ord}_q)$, with the derivative taken almost everywhere. The smooth-curve criterion says exactly that $b_q\geq0$, each $b_q$ is nonincreasing, and

\[
 \sum_q b_q(0+)\leq V,\qquad
 \varphi(t\operatorname{ord}_q)=a-\int_0^t b_q(s)\,ds.
\]

The endpoint value is the limit as $t\to\infty$, possibly $-\infty$. There are at most countably many nonconstant rays: their positive initial slopes have finite sum. Sums over $q$ mean the supremum over finite subsets. Define

\[
 D(\varphi)=\sum_q\int_0^\infty b_q(t)^2\,dt\in[0,\infty].
\]

For a PL potential, the raw model measure is $V\delta_o+\Delta\varphi$, by [BJ, Example 3.19] on $Y$ and the model normalization projection formula on $X$. Integration by parts on its finitely many rays yields

\[
 E_X(\varphi)=\frac1{2V}\int\varphi\,d(\mathrm{MA}^{\rm raw}_{X,L}(\varphi)+V\delta_o)
 =a-\frac{D(\varphi)}{2V}.
\]

This is the normalized intersection energy. The boundary term at infinity is zero for PL functions, which are eventually constant; the terms at the root are included in $V\delta_o$. Thus no conductor boundary term is omitted.

To extend the formula, approximate the functions $b_q$ from below by nonincreasing step functions with rational heights and rational breakpoints, supported on finite intervals of finitely many rays. Dyadic right-endpoint sampling, dyadic rounding down and increasing truncation give a directed increasing family of such approximants, converging almost everywhere to $b_q$. Their initial-slope sums do not exceed $V$. Their integrals give decreasing PL psh potentials with root value $a$, converging to $\varphi$, including at the endpoints. Real root constants are allowed in PL over $\mathbb R$. Monotone convergence of the squared slopes and the defining decreasing extension of energy give

\[
 E_X(\varphi)=a-\frac{D(\varphi)}{2V},\qquad
 \varphi\in\mathcal E^1(X,L)\Longleftrightarrow D(\varphi)<\infty.
\]

These are also approximations by the actual ample model cone. For clarity, if $f_n\downarrow\varphi$ is a countable PL approximation, choose $h_n\in H_{\rm dom}(X,L)$ with $\|h_n-f_n\|_\infty\leq\epsilon_n=2^{-n}$. Set $c_n=\sum_{j\geq n}(\epsilon_j+\epsilon_{j+1})=3\epsilon_n$. Then $h_n+c_n$ decreases, because $c_n-c_{n+1}=\epsilon_n+\epsilon_{n+1}$, and converges to $\varphi$. The constants are rational, as required for $H_{\rm dom}$. If $D(\varphi)<\infty$, the approximant energies decrease to a finite limit. Alternatively the source permits decreasing nets directly.

For $D(\varphi)<\infty$, every $b_q$ has limit zero at infinity. Define the operator explicitly by

\[
 \mathrm{MA}^{\rm raw}_{X,L}(\varphi)
 =\left(V-\sum_q b_q(0+)\right)\delta_o
   +\sum_q(-db_q)\big|_{(0,\infty)}.
\]

Here $-db_q$ is the positive Stieltjes measure on the $q$-ray, transported to actual $X$; right-continuous representatives are understood. Each ray measure has mass $b_q(0+)$, so the raw mass is exactly $V$. There is no mass at any closed-point endpoint. The countable sum of finite ray measures and the root atom is Radon. Finite partial sums are Radon and their total-variation tails tend to zero.

This is the genuine MA operator, not an independently named inverse measure. For every actual-X PL psh approximant it is the intersection measure, by model pullback and projection. Decreasing extension agrees with the unique operator of [BJ, Theorem 7.18 and Proposition 7.19(iv)]; continuous test functions determine the measures, since actual-X PL functions are dense. The same approximation proves $E_X(\varphi)=E_Y(\nu^*\varphi)$, first by projection for model pairings and then by their decreasing continuity. Consequently pullback identifies the finite-energy spaces, their weak-plus-energy strong topologies, and the genuine operators by

\[
 \mathrm{MA}^{\rm raw}_{X,L}(\varphi)=
 \nu_*\mathrm{MA}^{\rm raw}_{Y,A}(\nu^*\varphi).
\]

No source theorem is used to assume that arbitrary finite-energy potentials on $Y$ descend: Section 4 already proves that fact, and the energy identity now proves preservation of finite energy.

\section*{6. Exact measure condition and explicit existence and uniqueness}
Let $\mu$ be a positive Radon measure on actual $X^{\rm an}$ with raw mass $V$. Each ray minus $o$, including its endpoint, is open. A compact subset of their union meets finitely many rays, by the open cover by those rays. Inner regularity therefore shows that $\mu$ outside $o$ is concentrated on a countable union of rays. In particular this assertion does not assume that a measure is atomic, or supported on one skeleton. Write $\mu_q$ for its restriction to $(0,\infty]$ on the $q$-ray and set

\[
 \tau_q(t)=\mu_q((t,\infty]),\qquad
 D(\mu)=\sum_q\int_0^\infty\tau_q(t)^2\,dt.
\]

The endpoint is included in the tail in this definition. Endpoint mass makes $D(\mu)=\infty$. We will prove the exact identity

\[
 E_X^*(\mu/V)=\frac{D(\mu)}{2V},
\]

with the convention that both sides may be infinite.

For a finite-energy test potential with slope functions $b_q$ and root value $a$, Tonelli's theorem gives

\[
 -\int(\varphi-a)\,d\mu
 =\sum_q\int_0^\infty b_q(t)\tau_q(t)\,dt.
\]

If $D(\mu)<\infty$, Cauchy--Schwarz makes this finite, so every finite-energy potential is integrable for $\mu$. Substitution into the energy formula and completion of the square give

\[
 E_X(\varphi)-V^{-1}\int\varphi\,d\mu
 =V^{-1}\sum_q\int\left(b_q\tau_q-\tfrac12 b_q^2\right)dt
 \leq\frac{D(\mu)}{2V}.
\]

For the opposite inequality choose finitely many rays $F$ and $T<\infty$, and use the continuous psh test potential with root zero and slopes $b_q=\tau_q\mathbf1_{[0,T)}$ for $q\in F$, zero on other rays. Its slope budget is at most $V$, its energy is finite, and it is eventually constant on each active ray. It descends by Section 4. Its objective is exactly

\[
 \frac1{2V}\sum_{q\in F}\int_0^T\tau_q(t)^2\,dt.
\]

Taking the supremum over $F,T$ proves the identity even when $D(\mu)=\infty$. Uniform approximation of these continuous potentials by $H_{\rm dom}(X,L)$ preserves the objective. Conversely decreasing approximation of an arbitrary finite-energy potential and ordinary monotone convergence under the finite measure show that the defining supremum is unchanged by restricting to $H_{\rm dom}$. This proves both the finite-energy condition and its correct raw normalization.

Assume $D(\mu)<\infty$. Then $\mu$ has no endpoint mass. Define the root-normalized solution on actual $X$ by

\[
 \varphi(o)=0,\qquad
 \varphi(t\operatorname{ord}_q)
 =-\int_0^t\tau_q(s)\,ds
 =-\int_{(0,\infty)}\min(t,u)\,d\mu_q(u).
\]

Give each endpoint its decreasing limit. For finite $t$ the integral is finite, because $\tau_q\leq V$. The slopes are nonnegative, nonincreasing and have total initial value $\sum_q\mu_q((0,\infty))\leq V$. The finite-$F,T$ test potentials above decrease to this function. Each is psh on actual $X$, so decreasing closure proves the solution is psh, including usc at $o$ and all endpoints. Its energy is $-D(\mu)/(2V)>-\infty$. The Stieltjes formula gives $-d\tau_q=\mu_q$ on $(0,\infty)$, while the root coefficient is $\mu(\{o\})$. Hence its genuine raw MA measure is exactly $\mu$. This proves existence without a smooth-space solution being assumed to descend.

Conversely any finite-energy solution has no endpoint mass and its slope satisfies $-db_q=\mu_q$, $b_q(\infty)=0$. Thus $b_q(t)=\tau_q(t)$ at every continuity point, and everywhere for right-continuous representatives. Its energy formula forces $D(\mu)<\infty$. The slopes determine the entire potential up to its root value, including endpoint limits. This proves uniqueness modulo one constant and the necessity of the energy condition. It also gives $E_X^*(\mathrm{MA}^{\rm raw}(\varphi)/V)=\varphi(o)-E_X(\varphi)$.

\section*{7. Required analytic inputs, with each replacement checked}
Compactness. The weak topology is pointwise convergence on the common divisorial valuations. Section 4 identifies the two normalized psh cones as homeomorphic topological spaces. The cone on smooth connected $Y$ is compact by [BJ, Theorems 5.20 and 5.11]. Its image is therefore compact on actual singular $X$. The normalization also preserves $\sup\varphi=\varphi(o)$. We have not deleted smoothness from an SNC-model compactness proof; the proved cone identification supplies the replacement in this sector.

Envelopes. The homeomorphism of analytic spaces commutes with usc regularization, and the proved psh bijection identifies all competitors in an envelope. Thus for any bounded-above family and any continuous obstacle $f$,

\[
 \nu^*(\sup_i\varphi_i)^*=(\sup_i\nu^*\varphi_i)^*,\qquad
 \nu^*P_X(f)=P_Y(\nu^*f).
\]

The smooth-curve envelope property of [BJ, Theorem 5.20] now descends. For continuous $f$, $P_Y(\nu^*f)$ is continuous by Lemma 5.17, so $P_X(f)$ is continuous and psh. The same identities apply to bounded-above usc obstacles, with the identically negative-infinite case allowed. Neither normalization alone nor a presumed envelope property on $X$ was sufficient; the PL and psh descent proofs were necessary.

Energy differentiation. For $\varphi,\psi\in\mathcal E^1$, the quadratic energy formula gives, along their psh segment,

\[
 \left.\frac{d}{dt}\right|_{0+}E_X((1-t)\varphi+t\psi)
 =\frac1V\int(\psi-\varphi)\,d\mathrm{MA}^{\rm raw}_{X,L}(\varphi).
\]

The integrals are finite by the slope Cauchy--Schwarz inequality. This supplies segment differentiation directly. For the arbitrary continuous variations used in the variational method, [BJ, Theorem 8.5] gives the downward energy expansion $E_X^\downarrow(\varphi+\epsilon f)=E_X(\varphi)+\epsilon V^{-1}\int f\,d\mathrm{MA}^{\rm raw}(\varphi)+O(\epsilon^2)$ for actual-X PL $f$, uniformly in $\varphi$. It has no smoothness hypothesis. The now-proved envelope property identifies the left side with $E_X(P_X(\varphi+\epsilon f))$. Approximate a general continuous $f$ uniformly by actual-X PL $f_m$. The correct bound is $|f-f_m|\leq\epsilon_m$; monotonicity and constant translation of energy bound the difference of the perturbed energies by $|t|\epsilon_m$. Divide by $t$, use the PL derivative, and let $m\to\infty$. Consequently

\[
 \left.\frac{d}{dt}\right|_0 E_X(P_X(\varphi+tf))
 =V^{-1}\int f\,d\mathrm{MA}^{\rm raw}_{X,L}(\varphi)
 \quad (f\in C^0(X^{\rm an})).
\]

This also proves envelope orthogonality: for continuous $f$ put $u=P_X(f)$ and $g=f-u\geq0$. For $0\leq t\leq1$, every psh function below $u+tg\leq f$ is below $u$, while $u$ itself is a competitor. Hence $P_X(u+tg)=u$. Differentiation yields $\int(f-P_X(f))\,d\mathrm{MA}^{\rm raw}(P_X(f))=0$. Thus both differentiation and its envelope input are present; the explicit inverse construction is not being used in their place.

Comparison. [BJ, Theorem 7.40] applies to actual-X finite-energy potentials and its intersection operator, which Section 5 identifies with the constructed operator. It gives locality on $\{\varphi>\psi\}$ for $\max(\varphi,\psi)$. The maximum remains finite-energy by its slope/energy characterization, or by the cited finite-energy maximum property. Equal total mass and this locality, applied to $\max(\varphi,\psi-\epsilon)$, give

\[
 \int_{\{\varphi<\psi\}}\mathrm{MA}^{\rm raw}(\psi)
 \leq\int_{\{\varphi<\psi\}}\mathrm{MA}^{\rm raw}(\varphi).
\]

Indeed the partition into the two strict-inequality sets first bounds the left side over $\{\varphi<\psi-\epsilon\}$ by the right side over $\{\varphi\leq\psi-\epsilon\}$; both sets increase to $\{\varphi<\psi\}$ as $\epsilon\downarrow0$. Measures ignore the negative-infinite loci of finite-energy potentials; in this curve setting those possible loci are endpoints and have already been proved to have zero MA mass. This is the raw comparison inequality of Theorem 7.44. Uniqueness was also proved independently by the ray slopes.

Continuity and variational maximization. Theorem 7.19 gives decreasing continuity for any finite-energy approximating net, not only the selected step approximants. Proposition 12.7 gives strong continuity with strong potential topology equal to weak convergence plus energy convergence, and strong measure topology equal to weak measure convergence plus $E_X^*$ convergence. The exact psh and energy identifications preserve those topologies. Every finite-energy measure attains its Legendre supremum at the solution: the completed-square inequality in Section 6 is an equality at $b_q=\tau_q$. Thus the corresponding existence assertion also has the variational maximizer it requires. It is not deduced merely from completeness of a space of measures.

Reference dependence. The raw operator depends on the metric, not on the choice of reference used to name its potential. If $\rho$ is a continuous semipositive reference potential, the relative energy of the metric $\varphi$ is $E_{X,\rho}(\varphi-\rho)=E_X(\varphi)-E_X(\rho)$, by the intersection-energy cocycle formula and decreasing extension. Continuous references differ boundedly and identify the same finite-energy metrics. For raw mass $V$,

\[
 E_{X,\rho}^*(\mu/V)=E_X^*(\mu/V)-E_X(\rho)+V^{-1}\int\rho\,d\mu.
\]

The last two terms are finite, so the finite-energy condition is reference independent within this class. Multiplication of a metric by a constant shifts its potential and energy by the same additive constant and leaves both the raw measure and the dual objective unchanged. No singular reference type or nonample class is claimed here.

\section*{8. A concrete obstruction for a branched singular curve}
Let $X$ instead be an integral projective curve with an ordinary node $p$, and let $q_1,q_2$ be the two distinct points above $p$ on its smooth normalization. Retain algebraically closed trivially valued characteristic-zero $k$, an ample $L$, and $V=\deg L$. Let $v=\operatorname{ord}_{q_1}$, viewed as a point of $X^{\rm an}$, and put $\mu=V\delta_v$. This has the required raw mass.

It has finite dual energy on actual $X$. Any actual-X psh potential pulls back to a psh potential on $Y$; the energy is preserved by the model normalization projection formula and decreasing extension [BJ, (7.5)], which do not assert surjectivity of pullback. Therefore

\[
 E_X^*(\delta_v)\leq E_Y^*(\delta_{\operatorname{ord}_{q_1}})=V/2<\infty.
\]

The last equality is Section 6 on the smooth curve, with a raw mass $V$ atom at $t=1$. No arbitrary restriction of a measure has been normalized here.

Suppose $\mathrm{MA}^{\rm raw}_{X,L}(\varphi)=\mu$ for an actual-X finite-energy potential. Its pullback is finite-energy on $Y$. Formula (7.18) identifies the pushforward of its genuine MA with $\mu$. The analytic normalization is injective on finite generic-support ray points, and finite-energy MA has no endpoint mass. Thus the upstairs measure must be precisely $V\delta_{\operatorname{ord}_{q_1}}$. The unique normalized upstairs solution has slope $V$ on the $q_1$-ray up to $t=1$, zero afterwards, and is constant on the $q_2$-ray. Its two endpoint values are consequently $a-V$ and $a$, for some root value $a$. But both endpoints map to the single point $p$ of $X^{\rm an}$; a pulled-back function must have the same value at them. This contradiction proves that the finite-energy operator is not onto all finite-energy measures on this singular curve. It is the explicit obstruction recorded in [BJ, Example 12.11], now with raw mass and descent checked.

Thus a theorem for all singular varieties with only mass and the standard finite-energy condition is false in this framework. A broader branched theory would need a revised range condition, or a changed notion of potentials with an explicitly justified conductor theory. On a disconnected variety, each connected component has its own fixed polarization mass and its own additive constant; one global mass and uniqueness modulo a single constant are insufficient. The positive theorem avoids both issues by integrality and unibranchness.

\section*{9. Exact remaining original obligations}
The following gaps prevent candidate=true. First, no proof is given of the envelope property for an arbitrary higher-dimensional integral unibranch projective variety under retained characteristic-zero hypotheses. In the cached trivial-field source it is Conjecture 5.14, and Theorem 12.8 makes it precisely the missing existence input even though the operator, comparison and downward differentiation already exist there. This source-level statement is not a claim that no later theorem exists outside the authorized caches.

Second, the isolated-ray PL spanning argument is one-dimensional. For a higher-dimensional normalization, the needed cone descent and envelope identity have not been proved. Topological homeomorphism alone does not descend semipositivity, model approximation, energy or variational derivatives. Pulling an ample class back to a positive-dimensional resolution gives a generally nonample class, so a resolution cannot be substituted into an ample-class theorem. No such substitution was made above.

Third, no singular-space finite-energy operator, compactness, envelope, differentiation and comparison package is established here over a nontrivially discretely valued complete field of residue characteristic zero, in arbitrary dimension or for arbitrary finite-energy support. The cached BFJ v1 retains both algebraizability and SNC-support hypotheses. Dropping any of them needs its own exact later input and cannot follow from that v1. The curve theorem above does not dispose of this retained-field residual scope.

Fourth, for non-unibranch varieties the node counterexample requires an explicit corrected admissible-measure range and conductor descent theorem. This draft proves the obstruction and gives no general corrected range. Finally, irreducibility/connectedness and componentwise mass constraints must remain explicit in any wider theorem. These limitations keep the whole original atom open.

\section*{10. Reproducibility and evidence}
The files input-original.json and input.json are byte-identical; assignment-original.txt is the exact assignment string parsed from input.json, and task-original.json is its task object. Run \texttt{python3 provenance.py} to reproduce the selected literal source extracts and their application checks, using only the ten authorized caches. Run \texttt{python3 make\_paper.py} to produce the standalone P16.1-paper.tex from this complete argument. Run \texttt{python3 build\_and\_check.py} to build twice with the existing TeX installation, preserve both actual diagnostic streams and TeX logs, verify retained-input equality, validate the exact status.json key set, and execute the small exact-rational curve checks. Environment information is in environment.json. The checks are symbolic consistency examples and artifact checks, not an exhaustive result or a verification of the theorem.

The proof file and its standalone TeX/PDF contain the complete partial argument and all remaining original obligations. status.json records the partial result with exactly the required five keys. No other campaigns, generated notes, private payloads, external messages, installations, account changes, or heavy computation were used.

\end{document}
